\documentclass[12pt,a4paper]{article}
\usepackage[utf8]{vietnam}
\usepackage{amsmath,amssymb}
\usepackage{tabularx}
\usepackage{tikz}
\usetikzlibrary{calc,angles,quotes,arrows.meta,patterns}
\usepackage{xcolor}
\usepackage{enumitem}
\usepackage{fancyhdr}
\usepackage{graphicx}
\usepackage[top=0.9cm, bottom=0.9cm, left=1.2cm, right=1.2cm]{geometry}
\usepackage[dethi]{ex_test}
\SAOPTN{widthSA=3.2,heightSA=0.62}

\setlength{\headheight}{15pt}
\setlength{\parskip}{0pt}

\pagestyle{fancy}
\fancyhf{}
\rhead{\textit{Môn Toán 9 -- Đề số 06}}
\lhead{\textit{Bộ đề ôn tập Giữa học kì I}}
\cfoot{\thepage}

\begin{document}

% --- KHUNG TIÊU ĐỀ ĐỀ THI ---
\noindent
\begin{tabularx}{\linewidth}{X r}
\textbf{TRƯỜNG THCS: \dotfill} & \textbf{ĐỀ KIỂM TRA ĐỊNH KỲ GIỮA HỌC KỲ I} \\
\textbf{Họ và tên: \dotfill} & \textbf{NĂM HỌC 2026 -- 2027} \\
\textbf{Lớp: \dotfill \quad SBD: \dotfill} & \textbf{Môn: TOÁN 9} \\
\hline
\end{tabularx}
\vspace{2mm}

\begin{center}
  \textbf{\large ĐỀ ÔN TẬP GIỮA HỌC KỲ I -- ĐỀ SỐ 06}\\
  \textit{Thời gian làm bài: 90 phút (Không kể thời gian phát đề)}
\end{center}
\vspace{1mm}

% --- PHẦN I ---
\noindent\textbf{\large PHẦN I. CÂU TRẮC NGHIỆM NHIỀU PHƯƠNG ÁN LỰA CHỌN (3,0 điểm)}

\textit{Thí sinh trả lời từ câu 1 đến câu 12. Mỗi câu hỏi thí sinh chỉ chọn một phương án duy nhất.}

\vspace{1mm}
\begin{ex}
Phương trình nào dưới đây là phương trình bậc nhất hai ẩn?
\choice
{$2x^2 + 5y = 1$}
{\True $3x - 5y = 1$}
{$\frac{1}{x} + 2y = 1$}
{$4xy = 1$}
\end{ex}

\begin{ex}
Đường thẳng biểu diễn tất cả các nghiệm của phương trình $3x - 12 = 0$ là đường thẳng:
\choice
{Vuông góc với trục tung}
{Đi qua gốc tọa độ}
{Song song với trục hoành}
{\True Đi qua điểm $A(4; 5)$}
\end{ex}

\begin{ex}
Biết rằng $(x_0; y_0)$ là nghiệm của hệ phương trình $\begin{cases} x + 3y = 7 \\ 2x - y = 0. \end{cases}$ Tính $2y_0 - x_0$:
\choice
{$1$}
{$-2$}
{\True $3$}
{$-1$}
\end{ex}

\begin{ex}
Giải hệ phương trình $\begin{cases} x + y = 5 \\ x - y = 1 \end{cases}$ ta được nghiệm là:
\choice
{\True $(x; y) = (3; 2)$}
{$(x; y) = (2; 3)$}
{$(x; y) = (4; 1)$}
{$(x; y) = (1; 4)$}
\end{ex}

\begin{ex}
Phương trình $(x - 3)(x + 8) = 0$ có tập nghiệm là:
\choice
{$x = 3; x = 8$}
{$x = -3; x = 8$}
{$x = -3; x = -8$}
{\True $x = 3; x = -8$}
\end{ex}

\begin{ex}
Cho góc nhọn $B$ trong tam giác vuông. Tỉ số giữa cạnh kề và cạnh huyền là:
\choice
{$\sin B$}
{\True $\cos B$}
{$\tan B$}
{$\cot B$}
\end{ex}

\begin{ex}
Kết quả của phép tính $\sqrt{-13} : \sqrt{-468}$ (hiểu là $\sqrt{\frac{13}{468}}$) là:
\choice
{\True $\frac{1}{6}$}
{$-6$}
{$6$}
{$-\frac{1}{6}$}
\end{ex}

\begin{ex}
Đưa thừa số $-5a\sqrt{b}$ (với $a \ge 0, b \ge 0$) vào trong dấu căn ta được:
\choice
{\True $-\sqrt{25a^2 b}$}
{$\sqrt{25a^2 b}$}
{$-\sqrt{5a^2 b}$}
{$\sqrt{5a^2 b}$}
\end{ex}

\begin{ex}
Đưa thừa số $\sqrt{75x^2 y}$ (với $x < 0, y \ge 0$) ra ngoài dấu căn ta được:
\choice
{$5x\sqrt{3y}$}
{\True $-5x\sqrt{3y}$}
{$-75x\sqrt{y}$}
{$-25x\sqrt{3y}$}
\end{ex}

\begin{ex}
Cho $\alpha$ là góc nhọn bất kì có $\tan \alpha = \frac{1}{7}$, khi đó:
\choice
{$\cot \alpha = \frac{1}{7}$}
{$\cot \alpha = -\frac{1}{7}$}
{$\cot \alpha = -7$}
{\True $\cot \alpha = 7$}
\end{ex}

\begin{ex}
Chọn khẳng định đúng nhất trong các khẳng định sau:
\choice
{$\cos 77^\circ = \sin 23^\circ$}
{\True $\sin 37^\circ = \cos 53^\circ$}
{$\cot 53^\circ = \tan 47^\circ$}
{$\tan 37^\circ = \tan 53^\circ$}
\end{ex}

\begin{ex}
Trong một tam giác vuông, mỗi cạnh góc vuông bằng:
\choice
{Cạnh huyền nhân với tang góc đối hoặc côtang góc kề}
{Cạnh huyền nhân với côtang góc đối hoặc tang góc kề}
{\True Cạnh huyền nhân với sin góc đối hoặc côsin góc kề}
{Cạnh huyền nhân với côsin góc đối hoặc sin góc kề}
\end{ex}


\vspace{1mm}
\noindent\textbf{BẢNG TRẢ LỜI TRẮC NGHIỆM PHẦN I:}
\begin{center}
\begin{tabular}{|c|c|c|c|c|c|c|c|c|c|c|c|c|}
\hline
\textbf{Câu} & \textbf{1} & \textbf{2} & \textbf{3} & \textbf{4} & \textbf{5} & \textbf{6} & \textbf{7} & \textbf{8} & \textbf{9} & \textbf{10} & \textbf{11} & \textbf{12} \\
\hline
\textbf{Chọn} & \rule{0pt}{0.65cm} & & & & & & & & & & & \\
\hline
\end{tabular}
\end{center}
\vspace{2mm}

\setcounter{ex}{0}
\noindent\textbf{\large PHẦN II. CÂU TRẮC NGHIỆM ĐÚNG SAI (2,0 điểm)}

\textit{Thí sinh trả lời câu 1 và câu 2. Trong mỗi ý a), b), c), d) ở mỗi câu, thí sinh chọn đúng hoặc sai.}

\vspace{1mm}
\begin{ex}
Cho hệ phương trình $\begin{cases} 3x - 2y = 4 \\ x + 2y = 4. \end{cases}$ Xét tính Đúng/Sai của các khẳng định sau:
\choiceTF
{\True Hệ phương trình trên có hai ẩn số là $x$ và $y$}
{\True Cộng từng vế hai phương trình của hệ ta được phương trình $4x = 8$}
{Nghiệm của hệ phương trình là $(1; 2)$}
{\True Giá trị của biểu thức $T = x - y$ tại nghiệm của hệ bằng $1$}
\end{ex}

\begin{ex}
Cho tam giác $MNP$ vuông tại $M$, biết $MN = 9\text{ cm}$ và $MP = 12\text{ cm}$. Xét tính Đúng/Sai của các khẳng định sau:
\choiceTF
{\True Độ dài cạnh huyền $NP = 15\text{ cm}$}
{\True Tỉ số lượng giác $\sin N = \frac{4}{5}$}
{Số đo góc $N$ xấp xỉ bằng $37^\circ$}
{Tỉ số lượng giác $\tan P = \frac{4}{3}$}
\end{ex}


\noindent\textbf{BẢNG TRẢ LỜI PHẦN II:}
\begin{center}
\begin{tabular}{|c|c|c|c|c||c|c|c|c|c|}
\hline
\textbf{Câu 1} & \textbf{a)} & \textbf{b)} & \textbf{c)} & \textbf{d)} & \textbf{Câu 2} & \textbf{a)} & \textbf{b)} & \textbf{c)} & \textbf{d)} \\
\hline
\textbf{Đ/S} & \rule{0pt}{0.6cm} & & & & \textbf{Đ/S} & & & & \\
\hline
\end{tabular}
\end{center}
\vspace{2mm}

\setcounter{ex}{0}
\noindent\textbf{\large PHẦN III. CÂU HỎI TRẢ LỜI NGẮN (2,0 điểm)}

\textit{Thí sinh trả lời câu 1 và câu 2. Điền kết quả vào ô trống tương ứng.}

\vspace{1mm}
\begin{ex}
Giải hệ phương trình $\begin{cases} -2x + 5y = 10 \\ 2x - y = 2 \end{cases}$ ta được $y$ bằng bao nhiêu? \shortans{$3$}
\end{ex}

\begin{ex}
Tam giác $ABC$ có đường cao $AH$, biết $AB = 5\text{ cm}$ và $\widehat{B} = 60^\circ$. Hãy tính độ dài $AH$ (làm tròn đến hai chữ số thập phân, đơn vị cm). \shortans{$4{,}33$}
\end{ex}


\setcounter{ex}{0}
\noindent\textbf{\large PHẦN IV. TỰ LUẬN (3,0 điểm)}

\textit{Thí sinh trình bày chi tiết lời giải vào giấy làm bài.}

\vspace{1mm}
\begin{ex} (1,0 điểm)
Giải các phương trình sau:
\begin{enumerate}[label=\textbf{\alph*)}, leftmargin=0.8cm, itemsep=2pt]
\item (0,5 điểm) $(3x - 2)(2x + 1) = 0$
\item (0,5 điểm) $\frac{x + 1}{x - 3} - \frac{x - 2}{x + 3} = \frac{2(x + 5)}{x^2 - 9}$
\end{enumerate}
\end{ex}
\vspace{0.5mm}

\begin{ex} (0,75 điểm)
Một mảnh vườn hình chữ nhật có chu vi $50\text{ m}$. Nếu chiều dài tăng thêm $5\text{ m}$ và chiều rộng giảm đi $5\text{ m}$ thì diện tích của mảnh vườn giảm đi $50\text{ m}^2$. Tính diện tích của mảnh vườn đó.
\end{ex}
\vspace{0.5mm}

\begin{ex} (0,75 điểm)
Người ta làm một con đường gồm ba đoạn thẳng $AB, BC, CD$ bao quanh một hồ nước (như hình vẽ bên). Biết $AB = 10\text{ m}$, $BC = 13\text{ m}$, $CD = 15\text{ m}$ và góc tạo bởi đường $AB$ với phương vuông góc bờ $BC$ bằng $70^\circ$. Tính khoảng cách thẳng $AD$ (làm tròn đến hàng phần mười của mét).
\begin{center}
\begin{tikzpicture}[scale=0.35, line join=round, line cap=round, >=stealth]
  \coordinate (A) at (0,9.4);
  \coordinate (B) at (3.42,0);
  \coordinate (C) at (13,0);
  \coordinate (D) at (13,-15);
  \draw[fill=cyan!15, draw=blue!50] (4, -5) ellipse (6 and 4);
  \node at (4, -5) [blue!80!black] {Hồ nước};
  \draw[line width=1.8pt, color=orange!80!black] (A) -- (B) node[midway, above left] {$10\text{ m}$} -- (C) node[midway, above] {$13\text{ m}$} -- (D) node[midway, right] {$15\text{ m}$};
  \draw[thick, dashed, red] (A) -- (D) node[midway, left] {$AD$};
  \filldraw[black] (A) circle (3pt) node[above left] {$A$};
  \filldraw[black] (B) circle (3pt) node[below] {$B$};
  \filldraw[black] (C) circle (3pt) node[above right] {$C$};
  \filldraw[black] (D) circle (3pt) node[below right] {$D$};
\end{tikzpicture}
\end{center}
\end{ex}
\vspace{0.5mm}

\begin{ex} (0,5 điểm)
Tìm tất cả các cặp số nguyên $(x; y)$ thỏa mãn phương trình: $2x^2 + y^2 - 2xy + 10x + 15 = 0$.
\end{ex}
\vspace{0.5mm}


\enlargethispage{1.5cm}
\begin{center}
\rule{4cm}{0.4pt}\\
\textit{--- HẾT ---}
\end{center}
\end{document}
